\documentclass[11pt,a4paper]{article}
\usepackage{fontspec,amsmath,amssymb,amsthm,microtype}
\setmainfont{lmroman10-regular.otf}[BoldFont=lmroman10-bold.otf,ItalicFont=lmroman10-italic.otf,BoldItalicFont=lmroman10-bolditalic.otf]
\usepackage{xcolor}
\usepackage[most]{tcolorbox}
\usepackage{mathrsfs}
\providecommand{\lt}{<}
\providecommand{\gt}{>}
\setmonofont{DejaVuSansMono.ttf}[BoldFont=DejaVuSansMono-Bold.ttf,ItalicFont=DejaVuSansMono-Oblique.ttf,BoldItalicFont=DejaVuSansMono-BoldOblique.ttf,Scale=0.86]
\newfontfamily\leanfallback{STIXTwoMath-Regular.otf}[Scale=0.86]
\definecolor{leanrule}{gray}{0.62}
\definecolor{leanback}{gray}{0.962}
\definecolor{leantitle}{gray}{0.22}
\definecolor{leanmuted}{gray}{0.38}
\newtcolorbox{leanbox}[1]{enhanced,breakable,colback=leanback,colframe=leanrule,
  boxrule=0.4pt,arc=0.6mm,left=2.4mm,right=2.4mm,top=1.8mm,bottom=1.8mm,
  before skip=0.8em,after skip=1em,fonttitle=\small\bfseries,coltitle=leantitle,
  colbacktitle=leanback,titlerule=0.3pt,title={#1}}
\newlength{\leancharwidth}
\newenvironment{leancode}{\par\smallskip\begingroup\ttfamily\footnotesize\raggedright
  \settowidth{\leancharwidth}{M}\setlength{\parskip}{0pt}\setlength{\parindent}{0pt}
  \frenchspacing}{\par\endgroup\smallskip}
\newcommand{\leanline}[2]{\hangindent=\dimexpr#1\leancharwidth+4\leancharwidth\relax
  \hangafter=1\noindent\hspace*{#1\leancharwidth}#2\par}
\newcommand{\leansource}[1]{\par{\footnotesize\color{leanmuted}#1}\par}
\newcommand{\leanlabel}[1]{\par\smallskip\noindent{\small\bfseries #1}\par}
\usepackage[margin=25mm]{geometry}
\usepackage[unicode,hidelinks]{hyperref}
\hypersetup{pdfauthor={LangevinLean Project}}
\setlength{\emergencystretch}{3em}
\setlength{\parskip}{0.4em}
\setlength{\parindent}{0pt}
\allowdisplaybreaks[2]
\raggedbottom
\title{\texorpdfstring{Reversal symmetry of actual anchored kernel powers}{Reversal symmetry of actual anchored kernel powers}}
\author{LangevinLean Project}
\date{}
\begin{document}\maketitle
Draft mathematical exposition; these notes may contain mistakes.


One-step detailed balance implies reversal of every actual finite-step pair \(\mu(dx)K^n(x,dy)\). The proof uses nonnegative observable correlations and actual kernel integration, so the original initial coordinate remains present throughout.


\section{\texorpdfstring{Iterating the actual detailed-balance correlation}{Iterating the actual detailed-balance correlation}}


\textbf{Claim.} With \((Kf)(x)=\int f(y)K(x,dy)\), one-step pair reversal gives

\[\int f\,Kg\,d\mu=\int Kf\,g\,d\mu,\qquad\int f\,K^ng\,d\mu=\int K^nf\,g\,d\mu\]

for measurable nonnegative observables.



\textbf{Proof.} Product-kernel integration expresses the first pairing as the actual two-time integral. Swap invariance exchanges the two coordinates. For the iterated identity, move one averaging operator to the other observable by one-step balance, then apply the induction hypothesis. Kernel integration preserves measurability at every iterate.

\begin{leanbox}{Lean: \texttt{average}, \texttt{average\_\allowbreak{}measurable}, \texttt{average\_\allowbreak{}iterate\_\allowbreak{}measurable} and 2 more}
\leansource{\texttt{average} \quad WassersteinMixing/\allowbreak{}SemigroupPowers.lean, lines 26–26}
\begin{leancode}
\leanline{0}{def\ average\ (κ\ :\ Kernel\ E\ E)\ (f\ :\ E\ →\ ℝ≥0∞)\ (x\ :\ E)\ :\ ℝ≥0∞\ :=\ ∫⁻\ y,\ f\ y\ ∂κ\ x}
\end{leancode}
\leansource{\texttt{average\_\allowbreak{}measurable} \quad WassersteinMixing/\allowbreak{}SemigroupPowers.lean, lines 29–30}
\begin{leancode}
\leanline{0}{theorem\ average\_measurable\ (κ\ :\ Kernel\ E\ E)\ (f\ :\ E\ →\ ℝ≥0∞)\ (hf\ :\ Measurable\ f)\ :}
\leanline{4}{Measurable\ (average\ κ\ f)}
\end{leancode}
\leansource{\texttt{average\_\allowbreak{}iterate\_\allowbreak{}measurable} \quad WassersteinMixing/\allowbreak{}SemigroupPowers.lean, lines 33–39}
\begin{leancode}
\leanline{0}{theorem\ average\_iterate\_measurable\ (κ\ :\ Kernel\ E\ E)\ (f\ :\ E\ →\ ℝ≥0∞)}
\leanline{4}{(hf\ :\ Measurable\ f)\ (n\ :\ ℕ)\ :\ Measurable\ ((average\ κ)\textasciicircum{}[n]\ f)}
\end{leancode}
\leansource{\texttt{average\_\allowbreak{}balance} \quad WassersteinMixing/\allowbreak{}SemigroupPowers.lean, lines 42–63}
\begin{leancode}
\leanline{0}{theorem\ average\_balance\ (μ\ :\ Measure\ E)\ [SFinite\ μ]\ (κ\ :\ Kernel\ E\ E)}
\leanline{4}{[IsSFiniteKernel\ κ]\ (hswap\ :\ (μ\ ⊗ₘ\ κ).map\ Prod.swap\ =\ μ\ ⊗ₘ\ κ)}
\leanline{4}{(f\ g\ :\ E\ →\ ℝ≥0∞)\ (hf\ :\ Measurable\ f)\ (hg\ :\ Measurable\ g)\ :}
\leanline{4}{(∫⁻\ x,\ f\ x\ *\ average\ κ\ g\ x\ ∂μ)\ =\ ∫⁻\ x,\ average\ κ\ f\ x\ *\ g\ x\ ∂μ}
\end{leancode}
\leansource{\texttt{average\_\allowbreak{}iterate\_\allowbreak{}balance} \quad WassersteinMixing/\allowbreak{}SemigroupPowers.lean, lines 67–81}
\begin{leancode}
\leanline{0}{theorem\ average\_iterate\_balance\ (μ\ :\ Measure\ E)\ [SFinite\ μ]\ (κ\ :\ Kernel\ E\ E)}
\leanline{4}{[IsSFiniteKernel\ κ]\ (hswap\ :\ (μ\ ⊗ₘ\ κ).map\ Prod.swap\ =\ μ\ ⊗ₘ\ κ)}
\leanline{4}{(f\ g\ :\ E\ →\ ℝ≥0∞)\ (hf\ :\ Measurable\ f)\ (hg\ :\ Measurable\ g)\ (n\ :\ ℕ)\ :}
\leanline{4}{(∫⁻\ x,\ f\ x\ *\ ((average\ κ)\textasciicircum{}[n]\ g)\ x\ ∂μ)\ =}
\leanline{6}{∫⁻\ x,\ ((average\ κ)\textasciicircum{}[n]\ f)\ x\ *\ g\ x\ ∂μ}
\end{leancode}
\leanlabel{How the Lean proof follows the argument}
\texttt{average} is actual kernel integration. Its measurability lemmas justify each iteration. \texttt{average\_\allowbreak{}balance} changes variables under the swap law and \texttt{average\_\allowbreak{}iterate\_\allowbreak{}balance} proves the iterated pairing by induction.
\end{leanbox}

\section{\texorpdfstring{Actual powers, rectangles and their marginals}{Actual powers, rectangles and their marginals}}


\textbf{Claim.} Define \(K^0\) as the identity kernel and use actual kernel composition for successors. The pair law \(\mu\otimes K^n\) is swap invariant and has both marginals \(\mu\).



\textbf{Proof.} Kernel integration identifies the averaging action of the power with the iterated averaging function. Pair rectangle masses are indicator correlations, so the proved iterated balance gives equality with their reversed rectangles. Equality on rectangles proves full measure reversal. The first marginal is \(\mu\) because the kernel is Markov; swapping gives the second.

\begin{leanbox}{Lean: \texttt{kernelPower}, \texttt{kernelPower\_\allowbreak{}isMarkov}, \texttt{kernelPower\_\allowbreak{}average} and 4 more}
\leansource{\texttt{kernelPower} \quad WassersteinMixing/\allowbreak{}SemigroupPowers.lean, lines 84–86}
\begin{leancode}
\leanline{0}{def\ kernelPower\ (κ\ :\ Kernel\ E\ E)\ :\ ℕ\ →\ Kernel\ E\ E}
\leanline{2}{|\ 0\ =>\ Kernel.id}
\leanline{2}{|\ n\ +\ 1\ =>\ kernelPower\ κ\ n\ ∘ₖ\ κ}
\end{leancode}
\leansource{\texttt{kernelPower\_\allowbreak{}isMarkov} \quad WassersteinMixing/\allowbreak{}SemigroupPowers.lean, lines 88–92}
\begin{leancode}
\leanline{0}{instance\ kernelPower\_isMarkov\ (κ\ :\ Kernel\ E\ E)\ [IsMarkovKernel\ κ]\ (n\ :\ ℕ)\ :}
\leanline{4}{IsMarkovKernel\ (kernelPower\ κ\ n)\ :=\ by}
\leanline{2}{induction\ n\ with}
\leanline{2}{|\ zero\ =>\ unfold\ kernelPower;\ infer\_instance}
\leanline{2}{|\ succ\ n\ ih\ =>\ unfold\ kernelPower;\ infer\_instance}
\end{leancode}
\leansource{\texttt{kernelPower\_\allowbreak{}average} \quad WassersteinMixing/\allowbreak{}SemigroupPowers.lean, lines 95–107}
\begin{leancode}
\leanline{0}{theorem\ kernelPower\_average\ (κ\ :\ Kernel\ E\ E)\ (f\ :\ E\ →\ ℝ≥0∞)\ (hf\ :\ Measurable\ f)}
\leanline{4}{(n\ :\ ℕ)\ (x\ :\ E)\ :\ average\ (kernelPower\ κ\ n)\ f\ x\ =\ ((average\ κ)\textasciicircum{}[n]\ f)\ x}
\end{leancode}
\leansource{\texttt{kernelPower\_\allowbreak{}rectangle} \quad WassersteinMixing/\allowbreak{}SemigroupPowers.lean, lines 114–125}
\begin{leancode}
\leanline{0}{theorem\ kernelPower\_rectangle\ (n\ :\ ℕ)\ (s\ t\ :\ Set\ E)}
\leanline{4}{(hs\ :\ MeasurableSet\ s)\ (ht\ :\ MeasurableSet\ t)\ :}
\leanline{4}{(μ\ ⊗ₘ\ kernelPower\ κ\ n)\ (s\ ×ˢ\ t)\ =}
\leanline{6}{∫⁻\ x,\ s.indicator\ (fun\ \_\ ↦\ (1\ :\ ℝ≥0∞))\ x\ *}
\leanline{8}{((average\ κ)\textasciicircum{}[n]\ (t.indicator\ (fun\ \_\ ↦\ (1\ :\ ℝ≥0∞))))\ x\ ∂μ}
\end{leancode}
\leansource{\texttt{kernelPower\_\allowbreak{}pair\_\allowbreak{}swap} \quad WassersteinMixing/\allowbreak{}SemigroupPowers.lean, lines 129–145}
\begin{leancode}
\leanline{0}{theorem\ kernelPower\_pair\_swap}
\leanline{4}{(hswap\ :\ (μ\ ⊗ₘ\ κ).map\ Prod.swap\ =\ μ\ ⊗ₘ\ κ)\ (n\ :\ ℕ)\ :}
\leanline{4}{(μ\ ⊗ₘ\ kernelPower\ κ\ n).map\ Prod.swap\ =\ μ\ ⊗ₘ\ kernelPower\ κ\ n}
\end{leancode}
\leansource{\texttt{kernelPower\_\allowbreak{}pair\_\allowbreak{}fst} \quad WassersteinMixing/\allowbreak{}SemigroupPowers.lean, lines 148–150}
\begin{leancode}
\leanline{0}{theorem\ kernelPower\_pair\_fst\ (n\ :\ ℕ)\ :}
\leanline{4}{(μ\ ⊗ₘ\ kernelPower\ κ\ n).map\ Prod.fst\ =\ μ}
\end{leancode}
\leansource{\texttt{kernelPower\_\allowbreak{}pair\_\allowbreak{}snd} \quad WassersteinMixing/\allowbreak{}SemigroupPowers.lean, lines 153–157}
\begin{leancode}
\leanline{0}{theorem\ kernelPower\_pair\_snd}
\leanline{4}{(hswap\ :\ (μ\ ⊗ₘ\ κ).map\ Prod.swap\ =\ μ\ ⊗ₘ\ κ)\ (n\ :\ ℕ)\ :}
\leanline{4}{(μ\ ⊗ₘ\ kernelPower\ κ\ n).map\ Prod.snd\ =\ μ}
\end{leancode}
\leanlabel{How the Lean proof follows the argument}
\texttt{kernel\allowbreak{}Power} is the literal recursive kernel composition. \texttt{kernel\allowbreak{}Power\_\allowbreak{}average} and \texttt{kernel\allowbreak{}Power\_\allowbreak{}rectangle} connect it with the actual correlations. \texttt{kernel\allowbreak{}Power\_\allowbreak{}pair\_\allowbreak{}swap} uses product-measure extensionality. Its marginal theorems derive the exact laws.
\end{leanbox}


Earlier results used here: \cite[Section~3.35.5]{LangevinLean2026ULASC}.

\section{\texorpdfstring{Applying the power construction to the proved MALA pair}{Applying the power construction to the proved MALA pair}}


\textbf{Construction.} Disintegrate the actual reversible Gibbs MALA pair to obtain a Markov kernel. Its initial-law product is exactly the original pair. The actual kernel-power pair is therefore a probability measure, is reversible, and has both Gibbs marginals.



\textbf{Proof.} The disintegration theorem identifies the original pair with its first marginal times its conditional kernel. Its first marginal is already proved to be Gibbs. Apply the generic power reversal and marginal theorems to this exact one-step identity.

\begin{leanbox}{Lean: \texttt{kernel}, \texttt{kernel\_\allowbreak{}isMarkov}, \texttt{kernel\_\allowbreak{}pair\_\allowbreak{}eq} and 6 more}
\leansource{\texttt{kernel} \quad WassersteinMixing/\allowbreak{}SemigroupPowers.lean, lines 169–173}
\begin{leancode}
\leanline{0}{def\ kernel\ (V\ :\ E\ →\ ℝ)\ (h\ :\ ℝ)\ (L\ :\ ℝ≥0)}
\leanline{4}{[IsProbabilityMeasure\ (ULA.gibbsMeasure\ V)]}
\leanline{4}{(hdiff\ :\ Differentiable\ ℝ\ V)\ (hLip\ :\ LipschitzWith\ L\ (gradient\ V))\ :\ Kernel\ E\ E\ :=\ by}
\leanline{2}{have\ :=\ jointLaw\_isProbability\ V\ h\ L\ hdiff\ hLip}
\leanline{2}{exact\ (jointLaw\ V\ h).condKernel}
\end{leancode}
\leansource{\texttt{kernel\_\allowbreak{}isMarkov} \quad WassersteinMixing/\allowbreak{}SemigroupPowers.lean, lines 175–181}
\begin{leancode}
\leanline{0}{instance\ kernel\_isMarkov\ (V\ :\ E\ →\ ℝ)\ (h\ :\ ℝ)\ (L\ :\ ℝ≥0)}
\leanline{4}{[IsProbabilityMeasure\ (ULA.gibbsMeasure\ V)]}
\leanline{4}{(hdiff\ :\ Differentiable\ ℝ\ V)\ (hLip\ :\ LipschitzWith\ L\ (gradient\ V))\ :}
\leanline{4}{IsMarkovKernel\ (kernel\ V\ h\ L\ hdiff\ hLip)\ :=\ by}
\leanline{2}{have\ :=\ jointLaw\_isProbability\ V\ h\ L\ hdiff\ hLip}
\leanline{2}{unfold\ kernel}
\leanline{2}{infer\_instance}
\end{leancode}
\leansource{\texttt{kernel\_\allowbreak{}pair\_\allowbreak{}eq} \quad WassersteinMixing/\allowbreak{}SemigroupPowers.lean, lines 184–193}
\begin{leancode}
\leanline{0}{theorem\ kernel\_pair\_eq\ (V\ :\ E\ →\ ℝ)\ (h\ :\ ℝ)\ (L\ :\ ℝ≥0)}
\leanline{4}{[IsProbabilityMeasure\ (ULA.gibbsMeasure\ V)]}
\leanline{4}{(hdiff\ :\ Differentiable\ ℝ\ V)\ (hLip\ :\ LipschitzWith\ L\ (gradient\ V))\ :}
\leanline{4}{ULA.gibbsMeasure\ V\ ⊗ₘ\ kernel\ V\ h\ L\ hdiff\ hLip\ =\ jointLaw\ V\ h}
\end{leancode}
\leansource{\texttt{kernel\_\allowbreak{}pair\_\allowbreak{}swap} \quad WassersteinMixing/\allowbreak{}SemigroupPowers.lean, lines 196–202}
\begin{leancode}
\leanline{0}{theorem\ kernel\_pair\_swap\ (V\ :\ E\ →\ ℝ)\ (h\ :\ ℝ)\ (L\ :\ ℝ≥0)}
\leanline{4}{[IsProbabilityMeasure\ (ULA.gibbsMeasure\ V)]\ (hh\ :\ 0\ <\ h)}
\leanline{4}{(hdiff\ :\ Differentiable\ ℝ\ V)\ (hLip\ :\ LipschitzWith\ L\ (gradient\ V))\ :}
\leanline{4}{(ULA.gibbsMeasure\ V\ ⊗ₘ\ kernel\ V\ h\ L\ hdiff\ hLip).map\ Prod.swap\ =}
\leanline{6}{ULA.gibbsMeasure\ V\ ⊗ₘ\ kernel\ V\ h\ L\ hdiff\ hLip}
\end{leancode}
\leansource{\texttt{powerJointLaw} \quad WassersteinMixing/\allowbreak{}SemigroupPowers.lean, lines 206–210}
\begin{leancode}
\leanline{0}{def\ powerJointLaw\ (V\ :\ E\ →\ ℝ)\ (h\ :\ ℝ)\ (L\ :\ ℝ≥0)}
\leanline{4}{[IsProbabilityMeasure\ (ULA.gibbsMeasure\ V)]}
\leanline{4}{(hdiff\ :\ Differentiable\ ℝ\ V)\ (hLip\ :\ LipschitzWith\ L\ (gradient\ V))\ (n\ :\ ℕ)\ :}
\leanline{4}{Measure\ (E\ ×\ E)\ :=}
\leanline{2}{ULA.gibbsMeasure\ V\ ⊗ₘ\ StationaryKernelPower.kernelPower\ (kernel\ V\ h\ L\ hdiff\ hLip)\ n}
\end{leancode}
\leansource{\texttt{powerJointLaw\_\allowbreak{}isProbability} \quad WassersteinMixing/\allowbreak{}SemigroupPowers.lean, lines 212–217}
\begin{leancode}
\leanline{0}{instance\ powerJointLaw\_isProbability\ (V\ :\ E\ →\ ℝ)\ (h\ :\ ℝ)\ (L\ :\ ℝ≥0)}
\leanline{4}{[IsProbabilityMeasure\ (ULA.gibbsMeasure\ V)]}
\leanline{4}{(hdiff\ :\ Differentiable\ ℝ\ V)\ (hLip\ :\ LipschitzWith\ L\ (gradient\ V))\ (n\ :\ ℕ)\ :}
\leanline{4}{IsProbabilityMeasure\ (powerJointLaw\ V\ h\ L\ hdiff\ hLip\ n)\ :=\ by}
\leanline{2}{unfold\ powerJointLaw}
\leanline{2}{infer\_instance}
\end{leancode}
\leansource{\texttt{powerJointLaw\_\allowbreak{}swap} \quad WassersteinMixing/\allowbreak{}SemigroupPowers.lean, lines 220–225}
\begin{leancode}
\leanline{0}{theorem\ powerJointLaw\_swap\ (V\ :\ E\ →\ ℝ)\ (h\ :\ ℝ)\ (L\ :\ ℝ≥0)}
\leanline{4}{[IsProbabilityMeasure\ (ULA.gibbsMeasure\ V)]\ (hh\ :\ 0\ <\ h)}
\leanline{4}{(hdiff\ :\ Differentiable\ ℝ\ V)\ (hLip\ :\ LipschitzWith\ L\ (gradient\ V))\ (n\ :\ ℕ)\ :}
\leanline{4}{(powerJointLaw\ V\ h\ L\ hdiff\ hLip\ n).map\ Prod.swap\ =\ powerJointLaw\ V\ h\ L\ hdiff\ hLip\ n}
\end{leancode}
\leansource{\texttt{powerJointLaw\_\allowbreak{}fst} \quad WassersteinMixing/\allowbreak{}SemigroupPowers.lean, lines 228–233}
\begin{leancode}
\leanline{0}{theorem\ powerJointLaw\_fst\ (V\ :\ E\ →\ ℝ)\ (h\ :\ ℝ)\ (L\ :\ ℝ≥0)}
\leanline{4}{[IsProbabilityMeasure\ (ULA.gibbsMeasure\ V)]}
\leanline{4}{(hdiff\ :\ Differentiable\ ℝ\ V)\ (hLip\ :\ LipschitzWith\ L\ (gradient\ V))\ (n\ :\ ℕ)\ :}
\leanline{4}{(powerJointLaw\ V\ h\ L\ hdiff\ hLip\ n).map\ Prod.fst\ =\ ULA.gibbsMeasure\ V}
\end{leancode}
\leansource{\texttt{powerJointLaw\_\allowbreak{}snd} \quad WassersteinMixing/\allowbreak{}SemigroupPowers.lean, lines 236–241}
\begin{leancode}
\leanline{0}{theorem\ powerJointLaw\_snd\ (V\ :\ E\ →\ ℝ)\ (h\ :\ ℝ)\ (L\ :\ ℝ≥0)}
\leanline{4}{[IsProbabilityMeasure\ (ULA.gibbsMeasure\ V)]\ (hh\ :\ 0\ <\ h)}
\leanline{4}{(hdiff\ :\ Differentiable\ ℝ\ V)\ (hLip\ :\ LipschitzWith\ L\ (gradient\ V))\ (n\ :\ ℕ)\ :}
\leanline{4}{(powerJointLaw\ V\ h\ L\ hdiff\ hLip\ n).map\ Prod.snd\ =\ ULA.gibbsMeasure\ V}
\end{leancode}
\leanlabel{How the Lean proof follows the argument}
\texttt{kernel} is the actual conditional kernel of \texttt{joint\allowbreak{}Law}, not a postulated transition. \texttt{kernel\_\allowbreak{}pair\_\allowbreak{}eq} applies disintegration. \texttt{power\allowbreak{}Joint\allowbreak{}Law} and its law theorems instantiate the proved generic kernel-power results.
\end{leanbox}


Earlier results used here: \cite[Section~3.26.3]{LangevinLean2026ULASC}.

\begin{thebibliography}{ULA-LSI}
\bibitem[ULA-SC]{LangevinLean2026ULASC}
LangevinLean Project. \emph{Unadjusted Langevin under Strong Convexity}. LL-01, 2026. Draft report, content revision d34ead06987c. \url{https://odakyildiz.github.io/LangevinLean/docs/pdf/UnadjustedLangevin/FullProof.pdf}.
\end{thebibliography}

\end{document}
